\documentclass[11pt]{article}
\usepackage[margin=1in]{geometry}
\usepackage[T1]{fontenc}
\usepackage{lmodern}
\usepackage{amsmath,amssymb,amsthm}
\usepackage{graphicx}
\usepackage[colorlinks=true,linkcolor=blue,citecolor=blue,urlcolor=blue]{hyperref}
\emergencystretch=3em

\newtheorem{theorem}{Theorem}
\newtheorem{remark}{Remark}
\newcommand{\R}{\mathbb R}
\newcommand{\Z}{\mathbb Z}
\newcommand{\Sop}{\mathbf S}
\newcommand{\pointchain}[1]{[\![#1]\!]}

\title{Target Translations Obstruct Gradient-Only Control\\
of the Topological Singular-Set Operator}
\author{Run \texttt{fa\_banach\_001}}
\date{2026-08-17}

\begin{document}
\maketitle

\noindent\textbf{Status.}
Candidate full counterexample, likely valid, to the question following
Theorem~3.1 in the source paper.

\section{Source question}
Canevari and Orlandi construct
\[
 \Sop:(L^\infty\cap W^{1,k-1})(\Omega,\R^m)
 \longrightarrow
 L^1\!\left(\R^m;\mathbb F_{d-k}
 (\Omega;\pi_{k-1}(\mathcal N))\right)=:Y.
\]
After proving a continuity estimate involving $\|u_1-u_0\|_{L^k}$,
they ask whether $\|\Sop(u_1)-\Sop(u_0)\|_Y$ can instead be bounded in
terms of Sobolev norms of $\nabla u_1-\nabla u_0$~\cite{CO}.

\begin{center}
\includegraphics[width=0.98\linewidth]{figures/open_question_crop.png}
\end{center}

As is standard for a continuity estimate, ``controlled by the gradient
difference'' means that the proposed right-hand side vanishes when all
those norms vanish.  The example below rules out every such estimate,
including ones whose coefficient also depends on fixed $L^\infty$ bounds
or on individual derivative norms.

\section{Proof intuition}
The derivative forgets additive constants in the target, but $\Sop$ does
not.  In the sphere case, $\Sop_y(u)$ records the signed preimages of the
point $y$.  Replacing $u$ by $u+c$ therefore translates the whole family in
the parameter $y$, although it leaves the derivative unchanged.  Choosing
$c$ large enough makes the two ranges disjoint.  A positive-volume family
of unit point chains remains in their difference, so its $Y$-norm cannot
vanish.

\section{Counterexample theorem}
\begin{theorem}
For every integer $k\ge2$, let
\[
 d=m=k,\qquad \Omega=B_1(0)\subset\R^k,
 \qquad \mathcal N=S^{k-1}\subset\R^k.
\]
Put $c=3e_1$ and define smooth maps
\[
 u_0(x)=x,\qquad u_1(x)=x+c.
\]
Then
\[
 \nabla u_1-\nabla u_0=0
\]
pointwise, but
\[
 \|\Sop(u_1)-\Sop(u_0)\|_Y
 \ge \frac14\,|B_{1/2}(0)|>0.
\]
Consequently the source question has a negative answer in general: no
estimate controlled only by Sobolev norms of
$\nabla u_1-\nabla u_0$, and vanishing when that difference is zero, can
hold for $\Sop$.
\end{theorem}

\begin{proof}
Both maps belong to $L^\infty\cap W^{1,k-1}$ (indeed, they are affine), and
their derivatives are both the identity matrix.

For $\mathcal N=S^{k-1}\subset\R^k$, the source chooses the exceptional
complex $\mathcal X=\{0\}$ and the radial retraction
$z\mapsto z/|z|$.  Its smooth-map construction says that
$\Sop_y(u)$ is the signed zero-dimensional preimage chain of $y$: each
regular preimage carries its local degree in
$\pi_{k-1}(S^{k-1})\simeq\Z$~\cite[Sections 1.2 and 3.2]{CO}.  Since the
two affine derivatives preserve orientation, there is one fixed
$\varepsilon\in\{1,-1\}$, depending only on the source's generator
convention, such that for almost every $y$
\[
 \Sop_y(u_0)=
 \begin{cases}
  \varepsilon\pointchain{y},&y\in B_1(0),\\
  0,&y\notin B_1(0),
 \end{cases}
 \qquad
 \Sop_y(u_1)=
 \begin{cases}
  \varepsilon\pointchain{y-c},&y\in c+B_1(0),\\
  0,&y\notin c+B_1(0).
 \end{cases}
\]
The possible sign will play no role.

Because $|c|=3>2$, the two range balls are disjoint.  In particular, for
almost every $y\in B_{1/2}(0)$,
\[
 \Sop_y(u_1)-\Sop_y(u_0)
 =-\varepsilon\pointchain{y}
 \quad\text{in }\mathbb F_0(\Omega;\Z).
\]
It remains only to keep the relative flat norm from losing this point
through the boundary of $\Omega$.

Define the compactly supported Lipschitz function
\[
 \varphi(x)=\min\!\left\{\frac14,
                  \left(\frac34-|x|\right)_+\right\}.
\]
Then $0\le\varphi\le1$, $\operatorname{Lip}(\varphi)\le1$,
$\operatorname{spt}\varphi\Subset B_1$, and
$\varphi=1/4$ on $B_{1/2}$.  By the source's characterization of the
relative flat norm, any admissible decomposition of
$\pointchain{y}$ has the form
\[
   \pointchain{y}=\partial P+Q
 \quad\text{modulo a chain supported in }\R^k\setminus\Omega,
\]
and is charged by
$\mathbf M(P\mathbin{\llcorner}\Omega)
 +\mathbf M(Q\mathbin{\llcorner}\Omega)$.
Testing against $\varphi$ (or smooth compactly supported approximations)
gives, for $y\in B_{1/2}$,
\begin{align*}
 \frac14
 &=\pointchain{y}(\varphi)
   =P(d\varphi)+Q(\varphi)\\
 &\le
 \mathbf M(P\mathbin{\llcorner}\Omega)
 +\mathbf M(Q\mathbin{\llcorner}\Omega).
\end{align*}
Taking the infimum yields
$\mathbb F_\Omega(\pointchain{y})\ge1/4$.  Therefore
\[
 \|\Sop(u_1)-\Sop(u_0)\|_Y
 =\int_{\R^k}\mathbb F_\Omega
   (\Sop_y(u_1)-\Sop_y(u_0))\,dy
 \ge\int_{B_{1/2}}\frac14\,dy
 =\frac14|B_{1/2}|>0.
\]
This contradicts every proposed gradient-difference estimate whose
right-hand side vanishes at zero.
\end{proof}

\section{Scope, repair, and verification}
The counterexample addresses the literal unrestricted question.  It does
not conflict with the source's proved estimate, whose right-hand side
contains $\|u_1-u_0\|_{L^k}$.  Nor does it rule out a gradient-based estimate
after fixing target-translation freedom, for example by prescribing a
common mean, a common trace, or one point value.  Another natural repair is
to compare the $y$-families only after compensating for the target
translation.

The signed-preimage calculation is direct from the source's smooth
construction.  The relative-norm argument uses exactly its quotient-norm
characterization.  The accompanying verification script checks the affine
and disjoint-range identities in several dimensions; the proof above is
valid for every $k\ge2$.  A separate audit is recorded in
\texttt{VERIFICATION.md}.

\section{Bounded novelty search}
A search on 2026-08-17 used the exact question sentence, the paper title and
authors, ``topological singular set gradient difference'', and the
companion-paper and survey literature.  It found the source, Part~II
\cite{CO2}, the authors' 2021 survey~\cite{survey}, and later applications of
the same singular-chain operator, but no work explicitly answering this
question or recording the target-translation obstruction.  This is a
bounded search, not a priority claim.  Novelty confidence is moderate;
mathematical-validity confidence is high.

\begin{thebibliography}{9}
\bibitem{CO}
G. Canevari and G. Orlandi,
\emph{Topological singular set of vector-valued maps, I: Applications to
manifold-constrained Sobolev and BV spaces}, arXiv:1712.10203v2 (2018),
especially Theorem~3.1 and Sections~1.2, 2.1, 3.2.

\bibitem{CO2}
G. Canevari and G. Orlandi,
\emph{Topological singular set of vector-valued maps, II:
$\Gamma$-convergence for Ginzburg--Landau type functionals},
Arch. Ration. Mech. Anal. 241 (2021), 117--193; arXiv:2003.01354.

\bibitem{survey}
G. Canevari and G. Orlandi,
\emph{Topological singularities for vector-valued Sobolev maps and
applications}, Ann. Fac. Sci. Toulouse Math. 30 (2021), 327--351.
\end{thebibliography}

\end{document}
